\documentclass[12pt,a4paper]{article}
\usepackage[utf8]{inputenc}
\usepackage[T1]{fontenc}
\usepackage[spanish]{babel}
\usepackage{geometry}
\usepackage{graphicx}
\usepackage{enumitem}
\usepackage{booktabs}
\usepackage{longtable}
\usepackage{array}
\usepackage{xcolor}
\usepackage{titlesec}
\usepackage{fancyhdr}
\usepackage{hyperref}
\usepackage{colortbl}
\usepackage{float}

\geometry{margin=2.5cm}

% Colores
\definecolor{primary}{HTML}{FF9900}
\definecolor{darkprimary}{HTML}{E47911}
\definecolor{darkbg}{HTML}{232F3E}
\definecolor{lightgray}{HTML}{F5F5F5}
\definecolor{textgray}{HTML}{565959}

% Estilos de sección
\titleformat{\section}{\Large\bfseries\color{darkbg}}{}{0em}{}[\titlerule]
\titleformat{\subsection}{\large\bfseries\color{darkprimary}}{}{0em}{}
\titleformat{\subsubsection}{\normalsize\bfseries\color{darkbg}}{}{0em}{}

% Header y Footer
\pagestyle{fancy}
\fancyhf{}
\fancyhead[L]{\small\color{textgray}Marketplace de Fiestas --- Sprint 1}
\fancyhead[R]{\small\color{textgray}Unicornios Diabéticos}
\fancyfoot[C]{\thepage}
\renewcommand{\headrulewidth}{0.4pt}

\begin{document}

% ============================================================
% PORTADA
% ============================================================
\begin{titlepage}
\centering
\vspace*{2cm}

{\Huge\bfseries\color{darkbg} Entregables Técnicos y de\\[6pt] Documentación}\\[1.5cm]
{\LARGE\color{darkprimary} Marketplace con IA}\\[2cm]

\begin{tabular}{ll}
\textbf{Nombre del equipo:} & Unicornios Diabéticos \\[8pt]
\textbf{Integrante:} & Rodríguez Femat Emilio Emanuel \\[8pt]
\textbf{Área:} & Equipo de Front \\[8pt]
\textbf{Sprint:} & Sprint 1 \\[8pt]
\textbf{Fecha:} & 05 de marzo de 2026 \\[8pt]
\textbf{Versión:} & 1.0 \\
\end{tabular}

\vspace{2cm}

{\large\bfseries Actividades realizadas:}\\[0.7cm]
\begin{minipage}{0.8\textwidth}
\begin{itemize}[leftmargin=*]
    \item Definición de la lista de módulos del sistema
    \item Diseño de wireframes iniciales en Figma
    \item Elaboración del Product Backlog inicial
    \item Descripción de escenarios alternos y excepcionales
    \item Construcción de la matriz de navegación
    \item Elaboración de RF y RNF (35 RF y 17 RNF)
\end{itemize}
\end{minipage}

\vfill
\end{titlepage}

% ============================================================
% TABLA DE CONTENIDO
% ============================================================
\tableofcontents
\newpage

% ############################################################
% PARTE I — ENTREGABLES TÉCNICOS
% ############################################################
\part*{Parte I --- Entregables Técnicos}
\addcontentsline{toc}{part}{Parte I --- Entregables Técnicos}

% ============================================================
\section{Lista de Módulos del Sistema}
% ============================================================

El marketplace se organiza en \textbf{8 módulos funcionales} que agrupan los 35 requerimientos funcionales identificados en el levantamiento de requerimientos.

\subsection{Módulo 1: Autenticación}
\textbf{Requerimientos:} RF-01, RF-02, RF-03, RF-04, RF-05\\
Gestiona el registro de nuevos usuarios (comprador o vendedor), el inicio de sesión mediante correo y contraseña, la asignación de roles y el cierre de sesión seguro desde cualquier vista.

\subsection{Módulo 2: Catálogo y Búsqueda}
\textbf{Requerimientos:} RF-06, RF-07, RF-08, RF-09, RF-10, RF-11, RF-12\\
Presenta el catálogo general de productos organizados por categoría, permite la navegación y filtrado por categorías, ordenamiento por precio y calificación, y ofrece un buscador inteligente con autocompletado, tolerancia a errores tipográficos y sugerencias de productos similares.

\subsection{Módulo 3: Detalle de Producto}
\textbf{Requerimientos:} RF-13, RF-14, RF-15, RF-16\\
Muestra la información completa de cada producto (descripción, precio, imágenes, stock, reseñas, datos del vendedor y paquetes promocionales). Permite al comprador calificar productos adquiridos (1--5 estrellas con comentario) y guardarlos en su lista de deseos.

\subsection{Módulo 4: Carrito y Pedidos}
\textbf{Requerimientos:} RF-17, RF-18, RF-19, RF-20, RF-21, RF-22\\
Permite al comprador agregar productos al carrito, modificar cantidades o eliminar artículos, visualizar el resumen con precios y subtotales, y confirmar la compra generando una nota de pedido simulada con folio, fecha, lista de productos y monto total, sin solicitar datos bancarios. El carrito opera bajo un contexto fijo (<<Fiestas>>).

\subsection{Módulo 5: Perfil de Usuario}
\textbf{Requerimientos:} RF-23, RF-24, RF-25\\
Permite al usuario editar su información personal y registrar direcciones de envío. El comprador puede consultar su historial de pedidos anteriores y visualizar/gestionar los productos guardados en su lista de deseos.

\subsection{Módulo 6: Gestión de Productos (Vendedor)}
\textbf{Requerimientos:} RF-26, RF-27, RF-28, RF-29\\
Permite al vendedor publicar productos indicando nombre, descripción, precio, imágenes, stock y categoría. Incluye la edición y baja de productos, el control de cantidades en inventario y la creación de paquetes promocionales agrupando 2 o más artículos con precio preferencial.

\subsection{Módulo 7: Gestión de Pedidos (Vendedor)}
\textbf{Requerimientos:} RF-30, RF-31\\
Permite al vendedor consultar las órdenes recibidas y actualizar su estado (En preparación, Enviado, Entregado). Incluye un resumen básico de ventas con el total de órdenes y la calificación promedio.

\subsection{Módulo 8: Administración}
\textbf{Requerimientos:} RF-32, RF-33, RF-34, RF-35\\
Permite al administrador consultar la lista completa de usuarios registrados, bloquear/desbloquear o eliminar cuentas, y revisar reportes o denuncias realizadas por los usuarios sobre otros usuarios o productos inapropiados.

\begin{center}
\textbf{Total: 8 módulos | 35 requerimientos funcionales}
\end{center}

\newpage

% ============================================================
\section{Wireframes Iniciales}
% ============================================================

A continuación se presentan los wireframes de cada pantalla del sistema, diseñados en Figma. Cada sección incluye la descripción funcional, los requerimientos cubiertos y los elementos de interfaz representados.

% --- 2.1 ---
\subsection{Registro de Usuario}
\textbf{RF-01, RF-02 | Módulo: Autenticación | Actor: Comprador / Vendedor}

Pantalla que permite al usuario crear una cuenta proporcionando nombre completo, correo electrónico y contraseña, seleccionando su rol (comprador o vendedor). El sistema valida que el correo no esté previamente registrado.

\textbf{Elementos de la interfaz:}
\begin{itemize}[nosep]
    \item Logo del marketplace
    \item Campo: Nombre completo
    \item Campo: Correo electrónico
    \item Campo: Contraseña
    \item Selector de rol: Comprador / Vendedor / Administrador
    \item Botón: <<Iniciar Sesión>>

\end{itemize}

\begin{figure}[H]
\centering
\includegraphics[width=0.6\textwidth]{11.png} % <-- Cambia el nombre de la imagen aquí
\end{figure}



% --- 2.2 ---
\subsection{Inicio de Sesión}
\textbf{RF-03, RF-04, RF-05 | Módulo: Autenticación | Actor: Comprador / Vendedor / Administrador}

Pantalla de inicio de sesión mediante correo electrónico y contraseña. Redirige a la interfaz correspondiente según el rol del usuario. Permite cerrar sesión de forma segura desde cualquier vista.

\textbf{Elementos de la interfaz:}
\begin{itemize}[nosep]
    \item Logo del marketplace
    \item Campo: Nombre completo
    \item Campo: Correo electrónico
    \item Campo: Contraseña
    \item Campo: Confirmar contraseñas
    \item Botón: <<Iniciar sesión>>
    \item Enlace: <<Crear cuenta>>
\end{itemize}

\begin{figure}[H]
\centering
\includegraphics[width=0.35\textwidth]{2.png} % <-- Cambia el nombre de la imagen aquí
\end{figure}

% --- 2.3 ---
\subsection{Página Principal / Catálogo de Productos}
\textbf{RF-09, RF-10, RF-11, RF-12 | Módulo: Catálogo y Búsqueda | Actor: Comprador / Sistema}

Pantalla principal que muestra el catálogo general de productos organizados por categoría, una sección de productos destacados o mejor calificados, y opciones de filtrado y ordenamiento por precio y calificación.

\textbf{Elementos de la interfaz:}
\begin{itemize}[nosep]
    \item Barra de navegación: logo, búsqueda, cuenta, wishlist, carrito
    \item Banner/carrusel promocional
    \item Grid de categorías navegables
    \item Filtros: categoría, rango de precio, calificación
    \item Tarjetas de producto: imagen, nombre, precio
\end{itemize}

\begin{figure}[H]
\centering
\includegraphics[width=0.55\textwidth]{3.png} % <-- Cambia el nombre de la imagen aquí
\end{figure}



% --- 2.4 ---
\subsection{Búsqueda Inteligente}
\textbf{RF-06, RF-07, RF-08 | Módulo: Catálogo y Búsqueda | Actor: Comprador / Sistema}

Buscador inteligente con algoritmos de reconocimiento de patrones, autocompletado en tiempo real y tolerancia a errores tipográficos y sinónimos.

\textbf{Elementos de la interfaz:}
\begin{itemize}[nosep]
    \item Barra de búsqueda con ícono de lupa
    \item Dropdown de autocompletado con sugerencias
    \item Corrección de errores tipográficos (<<¿Quisiste decir...?>>)
    \item Sugerencias de productos similares
\end{itemize}

\begin{figure}[H]
\centering
\includegraphics[width=0.80\textwidth]{4.png} % <-- Cambia el nombre de la imagen aquí
\end{figure}


% --- 2.5 ---
\subsection{Detalle de Producto}
\textbf{RF-13, RF-14, RF-15, RF-16 | Módulo: Detalle de Producto | Actor: Comprador / Sistema}

Vista completa de un producto con toda su información, opciones de compra, paquetes promocionales del vendedor, reseñas de compradores y la posibilidad de calificar (solo si ya fue adquirido) o agregar a la lista de deseos.

\textbf{Elementos de la interfaz:}
\begin{itemize}[nosep]
    \item Galería de imágenes (principal + miniaturas)
    \item Nombre, precio, descripción del producto
    \item Indicador de stock disponible
    \item Rating en estrellas + conteo de reseñas
    \item Selector de cantidad
    \item Botón: <<Agregar al carrito>>
    \item Información del vendedor
    \item Sección de paquetes promocionales
    \item Sección de reseñas
    \item Formulario de calificación (estrellas + comentario)
\end{itemize}

\begin{figure}[H]
\centering
\includegraphics[width=0.55\textwidth]{5.png} % <-- Cambia el nombre de la imagen aquí
\end{figure}

\begin{figure}[H]
\centering
\includegraphics[width=0.65\textwidth]{6.png} % <-- Cambia el nombre de la imagen aquí
\end{figure}

% --- 2.6 ---
\subsection{Carrito de Compras}
\textbf{RF-17, RF-18, RF-19, RF-20 | Módulo: Carrito y Pedidos | Actor: Comprador / Sistema}

Muestra los artículos del carrito, permitiendo modificar cantidades o eliminar productos antes de confirmar. Opera bajo un contexto fijo de <<Fiestas>>.

\textbf{Elementos de la interfaz:}
\begin{itemize}[nosep]
    \item Lista de artículos: miniatura, nombre, precio unitario, selector de cantidad, subtotal, botón eliminar
    \item Indicador de contexto: <<Fiestas>>
    \item Resumen: subtotal, total
    \item Botón: <<Proceder al pago>>
\end{itemize}

\begin{figure}[H]
\centering
\includegraphics[width=0.55\textwidth]{7.png} % <-- Cambia el nombre de la imagen aquí
\end{figure}



% --- 2.7 ---
\subsection{Confirmación de Pedido (Checkout)}
\textbf{RF-21, RF-22 | Módulo: Carrito y Pedidos | Actor: Sistema}

Genera una nota de pedido simulada con folio, fecha, productos y monto total. No se solicitan datos bancarios. Tras confirmar, se muestra una pantalla de éxito.

\textbf{Elementos de la interfaz:}
\begin{itemize}[nosep]
    \item Resumen: folio, fecha, productos, cantidades, montos
    \item Dirección de envío seleccionada
    \item Monto total
    \item Botón: <<Confirmar pedido>>
    \item Pantalla de éxito con ícono de confirmación
    \item Botón: <<Seguir comprando>>
    \item NO se incluyen campos de tarjeta/datos bancarios
\end{itemize}

\begin{figure}[H]
\centering
\includegraphics[width=0.65\textwidth]{9.png} % <-- Cambia el nombre de la imagen aquí
\end{figure}

% --- 2.8 ---
\subsection{Perfil de Usuario / Mi Cuenta}
\textbf{RF-23, RF-05 | Módulo: Perfil de Usuario | Actor: Comprador / Vendedor}

Gestión del perfil para editar información personal y direcciones de envío. Punto de acceso al historial de pedidos y lista de deseos.

\textbf{Elementos de la interfaz:}
\begin{itemize}[nosep]
    \item Información personal editable (nombre, correo, teléfono)
    \item Sección de direcciones de envío (agregar, editar, eliminar)
    \item Navegación a: Historial de pedidos, Lista de deseos
    \item Botón: <<Cerrar sesión>>
\end{itemize}

\begin{figure}[H]
\centering
\includegraphics[width=0.65\textwidth]{10.png} % <-- Cambia el nombre de la imagen aquí
\end{figure}


% --- 2.9 ---
\subsection{Historial de Pedidos}
\textbf{RF-24 | Módulo: Perfil de Usuario | Actor: Comprador}

Listado de pedidos anteriores con fecha, productos comprados y monto total de cada pedido.

\textbf{Elementos de la interfaz:}
\begin{itemize}[nosep]
    \item Lista de pedidos: folio, fecha, miniaturas, total
    \item Opción para ver detalle de cada pedido
    \item Filtros por fecha
\end{itemize}

\begin{figure}[H]
\centering
\includegraphics[width=0.85\textwidth]{22.png} % <-- Cambia el nombre de la imagen aquí
\end{figure}






% --- 2.10 ---
\subsection{Lista de Deseos}
\textbf{RF-25 | Módulo: Perfil de Usuario | Actor: Comprador}

Productos guardados en la lista de deseos con opciones para gestionar.

\textbf{Elementos de la interfaz:}
\begin{itemize}[nosep]
    \item Grid de productos guardados
    \item Tarjetas: imagen, nombre, precio, disponibilidad
    \item Botón: <<Agregar al carrito>>
    \item Botón: <<Eliminar de la lista>>
    \item Estado vacío
\end{itemize}

\begin{figure}[H]
\centering
\includegraphics[width=0.55\textwidth]{12.png} % <-- Cambia el nombre de la imagen aquí
\end{figure}

% --- 2.11 ---
\subsection{Panel del Vendedor --- Gestión de Productos}
\textbf{RF-26, RF-27 | Módulo: Gestión de Productos (Vendedor) | Actor: Vendedor}

Panel para publicar nuevos productos y gestionar los existentes. Permite editar o dar de baja productos del catálogo.

\textbf{Elementos de la interfaz:}
\begin{itemize}[nosep]
    \item Tabla/lista de productos del vendedor
    \item Botón: <<Nuevo producto>>
    \item Formulario: nombre, descripción, precio, imágenes, stock, categoría
    \item Acciones: editar, dar de baja
    \item Navegación lateral del vendedor
\end{itemize}

\begin{figure}[H]
\centering
\includegraphics[width=0.85\textwidth]{13.png} % <-- Cambia el nombre de la imagen aquí
\end{figure}


% --- 2.12 ---
\subsection{Control de Inventario}
\textbf{RF-28 | Módulo: Gestión de Productos (Vendedor) | Actor: Vendedor}

Visualización y gestión de cantidades disponibles de cada producto.

\textbf{Elementos de la interfaz:}
\begin{itemize}[nosep]
    \item Tabla: nombre, stock actual, estado
    \item Campo editable de cantidad
    \item Indicadores: En stock / Bajo / Agotado
    \item Alertas de stock bajo
\end{itemize}

\begin{figure}[H]
\centering
\includegraphics[width=0.85\textwidth]{14.png} % <-- Cambia el nombre de la imagen aquí
\end{figure}


% --- 2.13 ---
\subsection{Paquetes Promocionales}
\textbf{RF-29 | Módulo: Gestión de Productos (Vendedor) | Actor: Vendedor}

Creación de paquetes promocionales agrupando 2 o más artículos con precio preferencial.

\textbf{Elementos de la interfaz:}
\begin{itemize}[nosep]
    \item Nombre del paquete
    \item Selector de productos (mínimo 2)
    \item Precios individuales
    \item Campo: precio preferencial del paquete
    \item Vista previa del paquete
    \item Botón: <<Crear paquete>>
\end{itemize}

\begin{figure}[H]
\centering
\includegraphics[width=0.85\textwidth]{15.png} % <-- Cambia el nombre de la imagen aquí
\end{figure}


% --- 2.14 ---
\subsection{Pedidos Recibidos del Vendedor}
\textbf{RF-30 | Módulo: Gestión de Pedidos (Vendedor) | Actor: Vendedor}

Órdenes recibidas de los compradores con opción de actualizar su estado.

\textbf{Elementos de la interfaz:}
\begin{itemize}[nosep]
    \item Lista/tabla de órdenes: folio, comprador, fecha, productos, total
    \item Selector de estado: En preparación / Enviado / Entregado
    \item Filtros por estado y fecha
    \item Detalle expandible por orden
\end{itemize}

\begin{figure}[H]
\centering
\includegraphics[width=0.85\textwidth]{16.png} % <-- Cambia el nombre de la imagen aquí
\end{figure}


% --- 2.15 ---
\subsection{Resumen de Ventas}
\textbf{RF-31 | Módulo: Gestión de Pedidos (Vendedor) | Actor: Vendedor}

Resumen básico de ventas con total de órdenes y calificación promedio.

\textbf{Elementos de la interfaz:}
\begin{itemize}[nosep]
    \item Tarjetas: total de órdenes, calificación promedio
    \item Gráfica básica de órdenes en el tiempo
    \item Lista de productos mejor calificados
\end{itemize}

\begin{figure}[H]
\centering
\includegraphics[width=0.85\textwidth]{23.png} % <-- Cambia el nombre de la imagen aquí
\end{figure}



% --- 2.16 ---
\subsection{Panel de Administración --- Gestión de Usuarios}
\textbf{RF-32, RF-33, RF-34 | Módulo: Administración | Actor: Administrador}

Lista completa de usuarios registrados con opciones para bloquear, desbloquear o eliminar cuentas.

\textbf{Elementos de la interfaz:}
\begin{itemize}[nosep]
    \item Tabla: nombre, correo, rol, fecha registro, estado
    \item Búsqueda de usuarios
    \item Filtros por rol y estado
    \item Acciones: ver perfil y eliminar
    \item Modal de confirmación para acciones destructivas
    \item Contador de usuarios registrados
\end{itemize}

\begin{figure}[H]
\centering
\includegraphics[width=0.85\textwidth]{31.png} % <-- Cambia el nombre de la imagen aquí
\end{figure}

% --- 2.17 ---
\subsection{Reportes y Denuncias}
\textbf{RF-35 | Módulo: Administración | Actor: Administrador}

Revisión de reportes o denuncias sobre usuarios o productos inapropiados.

\textbf{Elementos de la interfaz:}
\begin{itemize}[nosep]
    \item Lista de reportes: ID, denunciante, reportado, motivo, fecha, estado
    \item Vista de detalle del reporte
    \item Acciones: desestimar, advertir, bloquear, eliminar producto
    \item Filtros por tipo y estado
\end{itemize}

\begin{figure}[H]
\centering
\includegraphics[width=0.85\textwidth]{32.png} % <-- Cambia el nombre de la imagen aquí
\end{figure}



\newpage

% ############################################################
% PARTE II — ENTREGABLES DE DOCUMENTACIÓN
% ############################################################
\part*{Parte II --- Entregables de Documentación}
\addcontentsline{toc}{part}{Parte II --- Entregables de Documentación}

% ============================================================
\section{Product Backlog Inicial}
% ============================================================

El Product Backlog se organiza por prioridad, basado en el valor de negocio y las dependencias técnicas entre los módulos. Cada historia de usuario se asocia a su requerimiento funcional correspondiente.

\textbf{Prioridades:}
\begin{itemize}[nosep]
    \item \textbf{Alta} = Funcionalidad esencial para el MVP del marketplace.
    \item \textbf{Media} = Funcionalidad importante pero no bloqueante.
\end{itemize}

\textbf{Estimación en Story Points:} 1 = Muy simple | 2 = Simple | 3 = Moderado | 5 = Complejo | 8 = Muy complejo

% --- Sprint 1 ---
\subsection{Sprint 1 --- Autenticación y Estructura Base}
\begin{longtable}{|c|p{8cm}|c|c|c|}
\hline
\rowcolor{darkbg}\textcolor{white}{\textbf{ID}} & \textcolor{white}{\textbf{Historia de Usuario}} & \textcolor{white}{\textbf{RF}} & \textcolor{white}{\textbf{Prior.}} & \textcolor{white}{\textbf{Pts}} \\
\hline
HU-01 & Como usuario, quiero registrarme con mi nombre, correo y contraseña seleccionando mi rol para acceder a la plataforma. & RF-01 & Alta & 3 \\
\hline
HU-02 & Como sistema, quiero validar que el correo no esté registrado previamente para evitar cuentas duplicadas. & RF-02 & Alta & 2 \\
\hline
HU-03 & Como usuario, quiero iniciar sesión con mi correo y contraseña para acceder a mi cuenta. & RF-03 & Alta & 2 \\
\hline
HU-04 & Como sistema, quiero mostrar la interfaz correspondiente según el rol del usuario autenticado. & RF-04 & Alta & 3 \\
\hline
HU-05 & Como usuario, quiero cerrar sesión de forma segura desde cualquier vista para proteger mi cuenta. & RF-05 & Alta & 1 \\
\hline
\end{longtable}

% --- Sprint 2 ---
\subsection{Sprint 2 --- Catálogo, Búsqueda y Navegación}
\begin{longtable}{|c|p{8cm}|c|c|c|}
\hline
\rowcolor{darkbg}\textcolor{white}{\textbf{ID}} & \textcolor{white}{\textbf{Historia de Usuario}} & \textcolor{white}{\textbf{RF}} & \textcolor{white}{\textbf{Prior.}} & \textcolor{white}{\textbf{Pts}} \\
\hline
HU-06 & Como comprador, quiero buscar productos usando un buscador inteligente que analice mis términos y los compare con nombres, descripciones y etiquetas. & RF-06 & Alta & 8 \\
\hline
HU-07 & Como comprador, quiero que el buscador me ofrezca autocompletado conforme escribo para encontrar productos más rápido. & RF-07 & Media & 5 \\
\hline
HU-08 & Como comprador, quiero que el buscador sugiera productos similares aunque escriba con errores tipográficos o sinónimos. & RF-08 & Media & 5 \\
\hline
HU-09 & Como comprador, quiero ver un catálogo general de productos organizados por categoría. & RF-09 & Alta & 3 \\
\hline
HU-10 & Como comprador, quiero filtrar productos por categorías predefinidas para encontrar lo que busco. & RF-10 & Alta & 3 \\
\hline
HU-11 & Como comprador, quiero ordenar los resultados por precio y calificación para comparar opciones. & RF-11 & Alta & 2 \\
\hline
HU-12 & Como comprador, quiero ver productos destacados o mejor calificados en la página principal. & RF-12 & Media & 3 \\
\hline
\end{longtable}

% --- Sprint 3 ---
\subsection{Sprint 3 --- Detalle de Producto, Calificaciones y Wishlist}
\begin{longtable}{|c|p{8cm}|c|c|c|}
\hline
\rowcolor{darkbg}\textcolor{white}{\textbf{ID}} & \textcolor{white}{\textbf{Historia de Usuario}} & \textcolor{white}{\textbf{RF}} & \textcolor{white}{\textbf{Prior.}} & \textcolor{white}{\textbf{Pts}} \\
\hline
HU-13 & Como comprador, quiero ver el detalle de un producto con descripción, precio, imágenes, stock, reseñas y datos del vendedor. & RF-13 & Alta & 5 \\
\hline
HU-14 & Como comprador, quiero ver los paquetes promocionales del vendedor como opción de compra. & RF-14 & Media & 3 \\
\hline
HU-15 & Como comprador, quiero calificar de 1 a 5 estrellas y comentar productos que ya haya comprado. & RF-15 & Media & 3 \\
\hline
HU-16 & Como comprador, quiero guardar productos en una lista de deseos para consultarlos después. & RF-16 & Media & 2 \\
\hline
\end{longtable}

% --- Sprint 4 ---
\subsection{Sprint 4 --- Carrito de Compras y Checkout}
\begin{longtable}{|c|p{8cm}|c|c|c|}
\hline
\rowcolor{darkbg}\textcolor{white}{\textbf{ID}} & \textcolor{white}{\textbf{Historia de Usuario}} & \textcolor{white}{\textbf{RF}} & \textcolor{white}{\textbf{Prior.}} & \textcolor{white}{\textbf{Pts}} \\
\hline
HU-17 & Como comprador, quiero agregar productos al carrito desde el detalle del producto. & RF-17 & Alta & 2 \\
\hline
HU-18 & Como sistema, quiero restringir el carrito al contexto fijo de <<Fiestas>>. & RF-18 & Alta & 1 \\
\hline
HU-19 & Como comprador, quiero ver el resumen de mi carrito con artículos, cantidades, precio y subtotal. & RF-19 & Alta & 3 \\
\hline
HU-20 & Como comprador, quiero modificar cantidades o eliminar productos del carrito antes de confirmar. & RF-20 & Alta & 2 \\
\hline
HU-21 & Como sistema, quiero generar una nota de pedido simulada con folio, fecha, productos y total al confirmar la compra, sin solicitar datos bancarios. & RF-21 & Alta & 5 \\
\hline
HU-22 & Como comprador, quiero ver una pantalla de confirmación con el resumen de mi pedido. & RF-22 & Alta & 2 \\
\hline
\end{longtable}

% --- Sprint 5 ---
\subsection{Sprint 5 --- Perfil, Historial y Lista de Deseos}
\begin{longtable}{|c|p{8cm}|c|c|c|}
\hline
\rowcolor{darkbg}\textcolor{white}{\textbf{ID}} & \textcolor{white}{\textbf{Historia de Usuario}} & \textcolor{white}{\textbf{RF}} & \textcolor{white}{\textbf{Prior.}} & \textcolor{white}{\textbf{Pts}} \\
\hline
HU-23 & Como usuario, quiero editar mi información personal y registrar direcciones de envío. & RF-23 & Alta & 3 \\
\hline
HU-24 & Como comprador, quiero consultar el historial de mis pedidos con fecha, productos y monto total. & RF-24 & Media & 3 \\
\hline
HU-25 & Como comprador, quiero visualizar y gestionar los productos de mi lista de deseos. & RF-25 & Media & 2 \\
\hline
\end{longtable}

% --- Sprint 6 ---
\subsection{Sprint 6 --- Panel del Vendedor: Productos e Inventario}
\begin{longtable}{|c|p{8cm}|c|c|c|}
\hline
\rowcolor{darkbg}\textcolor{white}{\textbf{ID}} & \textcolor{white}{\textbf{Historia de Usuario}} & \textcolor{white}{\textbf{RF}} & \textcolor{white}{\textbf{Prior.}} & \textcolor{white}{\textbf{Pts}} \\
\hline
HU-26 & Como vendedor, quiero publicar productos con nombre, descripción, precio, imágenes, stock y categoría. & RF-26 & Alta & 5 \\
\hline
HU-27 & Como vendedor, quiero editar o dar de baja productos existentes de mi catálogo. & RF-27 & Alta & 3 \\
\hline
HU-28 & Como vendedor, quiero visualizar y gestionar las cantidades de inventario de mis productos. & RF-28 & Alta & 3 \\
\hline
HU-29 & Como vendedor, quiero crear paquetes promocionales agrupando 2 o más artículos con precio preferencial. & RF-29 & Media & 5 \\
\hline
\end{longtable}

% --- Sprint 7 ---
\subsection{Sprint 7 --- Panel del Vendedor: Pedidos y Ventas}
\begin{longtable}{|c|p{8cm}|c|c|c|}
\hline
\rowcolor{darkbg}\textcolor{white}{\textbf{ID}} & \textcolor{white}{\textbf{Historia de Usuario}} & \textcolor{white}{\textbf{RF}} & \textcolor{white}{\textbf{Prior.}} & \textcolor{white}{\textbf{Pts}} \\
\hline
HU-30 & Como vendedor, quiero consultar las órdenes recibidas y actualizar su estado (En preparación, Enviado, Entregado). & RF-30 & Alta & 5 \\
\hline
HU-31 & Como vendedor, quiero ver un resumen de mis ventas con total de órdenes y calificación promedio. & RF-31 & Media & 3 \\
\hline
\end{longtable}

% --- Sprint 8 ---
\subsection{Sprint 8 --- Panel de Administración}
\begin{longtable}{|c|p{8cm}|c|c|c|}
\hline
\rowcolor{darkbg}\textcolor{white}{\textbf{ID}} & \textcolor{white}{\textbf{Historia de Usuario}} & \textcolor{white}{\textbf{RF}} & \textcolor{white}{\textbf{Prior.}} & \textcolor{white}{\textbf{Pts}} \\
\hline
HU-32 & Como administrador, quiero consultar la lista completa de usuarios registrados. & RF-32 & Alta & 3 \\
\hline
HU-33 & Como administrador, quiero bloquear o desbloquear cuentas de usuarios que incumplan las normas. & RF-33 & Alta & 3 \\
\hline
HU-34 & Como administrador, quiero eliminar cuentas de usuarios en casos graves. & RF-34 & Alta & 2 \\
\hline
HU-35 & Como administrador, quiero revisar reportes o denuncias sobre usuarios o productos inapropiados. & RF-35 & Media & 5 \\
\hline
\end{longtable}

% --- Resumen ---
\subsection{Resumen del Product Backlog}

\begin{center}
\begin{tabular}{ll}
\toprule
\textbf{Total de Historias de Usuario} & 35 \\
\textbf{Total de Story Points} & 113 \\
\midrule
\textbf{Prioridad Alta} & 22 historias (67 puntos) \\
\textbf{Prioridad Media} & 13 historias (46 puntos) \\
\bottomrule
\end{tabular}
\end{center}

\begin{center}
\begin{tabular}{lcc}
\toprule
\textbf{Sprint} & \textbf{HU} & \textbf{Puntos} \\
\midrule
Sprint 1 --- Autenticación & 5 & 11 \\
Sprint 2 --- Catálogo y Búsqueda & 7 & 29 \\
Sprint 3 --- Detalle y Calificaciones & 4 & 13 \\
Sprint 4 --- Carrito y Checkout & 6 & 15 \\
Sprint 5 --- Perfil e Historial & 3 & 8 \\
Sprint 6 --- Vendedor: Productos & 4 & 16 \\
Sprint 7 --- Vendedor: Pedidos & 2 & 8 \\
Sprint 8 --- Administración & 4 & 13 \\
\bottomrule
\end{tabular}
\end{center}

\newpage

% ============================================================
\section{Descripción de Posibles escenarios}
% ============================================================

A continuación se presentan los escenarios alternos y excepcionales de cada pantalla, describiendo qué sucede cuando el usuario realiza acciones inesperadas, ocurren errores o se presentan condiciones especiales (los cuales se pueden tener en cuenta posteriormente).

\subsection{Registro de Usuario (RF-01, RF-02)}
\begin{itemize}
    \item \textbf{¿Qué pasa si el usuario intenta registrarse con un correo ya existente?} $\rightarrow$ El sistema muestra un mensaje de error: <<Este correo ya está registrado. Intenta iniciar sesión o usa otro correo.>>
    \item \textbf{¿Qué pasa si el usuario deja campos obligatorios vacíos?} $\rightarrow$ El sistema resalta los campos faltantes en rojo y muestra: <<Por favor completa todos los campos requeridos.>>
    \item \textbf{¿Qué pasa si la contraseña no cumple los requisitos mínimos?} $\rightarrow$ El sistema indica: <<La contraseña debe tener al menos 8 caracteres.>>
    \item \textbf{¿Qué pasa si el usuario no selecciona un rol?} $\rightarrow$ El botón <<Crear cuenta>> permanece deshabilitado hasta que se seleccione un rol.
\end{itemize}

\subsection{Inicio de Sesión (RF-03, RF-04, RF-05)}
\begin{itemize}
    \item \textbf{¿Qué pasa si el usuario ingresa un correo o contraseña incorrectos?} $\rightarrow$ El sistema muestra: <<Correo o contraseña incorrectos. Intenta de nuevo.>>
    \item \textbf{¿Qué pasa si el usuario intenta iniciar sesión con una cuenta bloqueada?} $\rightarrow$ El sistema muestra: <<Tu cuenta ha sido bloqueada. Contacta al administrador.>>
    \item \textbf{¿Qué pasa si el usuario cierra sesión?} $\rightarrow$ Se redirige a la pantalla de login y se invalida la sesión activa.
\end{itemize}

\subsection{Página Principal / Catálogo (RF-09, RF-10, RF-11, RF-12)}
\begin{itemize}
    \item \textbf{¿Qué pasa si no hay productos disponibles en una categoría?} $\rightarrow$ Se muestra: <<No hay productos disponibles en esta categoría por el momento.>>
    \item \textbf{¿Qué pasa si el usuario aplica filtros y no hay resultados?} $\rightarrow$ Se muestra: <<No se encontraron productos con los filtros seleccionados.>>
    \item \textbf{¿Qué pasa si el usuario no está autenticado?} $\rightarrow$ Puede navegar y ver productos, pero no puede agregar al carrito ni a la lista de deseos.
\end{itemize}

\subsection{Búsqueda Inteligente (RF-06, RF-07, RF-08)}
\begin{itemize}
    \item \textbf{¿Qué pasa si el término no coincide con ningún producto?} $\rightarrow$ El sistema muestra: <<No se encontraron resultados. ¿Quisiste decir...?>> con sugerencias.
    \item \textbf{¿Qué pasa si el usuario escribe con errores tipográficos?} $\rightarrow$ El buscador aplica tolerancia a errores y muestra resultados aproximados.
    \item \textbf{¿Qué pasa si el usuario envía la búsqueda vacía?} $\rightarrow$ Se muestran los productos destacados o se redirige al catálogo general.
\end{itemize}

\subsection{Detalle de Producto (RF-13, RF-14, RF-15, RF-16)}
\begin{itemize}
    \item \textbf{¿Qué pasa si el producto no tiene stock?} $\rightarrow$ Se muestra <<Agotado>> y el botón de agregar al carrito se deshabilita.
    \item \textbf{¿Qué pasa si intenta calificar un producto que no ha comprado?} $\rightarrow$ El formulario aparece deshabilitado con el mensaje: <<Solo puedes calificar productos que hayas comprado.>>
    \item \textbf{¿Qué pasa si intenta agregar a deseos sin estar autenticado?} $\rightarrow$ Se redirige a la pantalla de inicio de sesión.
    \item \textbf{¿Qué pasa si el producto no tiene reseñas?} $\rightarrow$ Se muestra: <<Este producto aún no tiene reseñas. ¡Sé el primero en opinar!>>
\end{itemize}

\subsection{Carrito de Compras (RF-17, RF-18, RF-19, RF-20)}
\begin{itemize}
    \item \textbf{¿Qué pasa si intenta agregar más cantidad que el stock?} $\rightarrow$ El sistema limita la cantidad al stock disponible.
    \item \textbf{¿Qué pasa si el carrito está vacío e intenta ir al checkout?} $\rightarrow$ El botón <<Proceder al pago>> está deshabilitado.
    \item \textbf{¿Qué pasa si un producto se agota mientras navega?} $\rightarrow$ Se notifica y se marca como no disponible en el carrito.
    \item \textbf{¿Qué pasa si elimina todos los productos?} $\rightarrow$ Se muestra el estado vacío con botón para seguir comprando.
\end{itemize}

\subsection{Confirmación de Pedido (RF-21, RF-22)}
\begin{itemize}
    \item \textbf{¿Qué pasa si no tiene dirección de envío?} $\rightarrow$ Se le solicita registrar una antes de confirmar.
    \item \textbf{¿Qué pasa si no tiene internet al confirmar?} $\rightarrow$ Se muestra: <<No se pudo procesar tu pedido. Verifica tu conexión.>>
    \item \textbf{¿Qué pasa si cierra la página después de confirmar?} $\rightarrow$ El pedido ya fue registrado y puede consultarlo en su historial.
\end{itemize}

\subsection{Perfil de Usuario (RF-23, RF-05)}
\begin{itemize}
    \item \textbf{¿Qué pasa si guarda información incompleta?} $\rightarrow$ Se resaltan los campos obligatorios faltantes.
    \item \textbf{¿Qué pasa si elimina su única dirección?} $\rightarrow$ Se permite pero se advierte que necesitará una para comprar.
    \item \textbf{¿Qué pasa si pone mal su contraseña al cambiar datos?} $\rightarrow$ Se solicita verificación de contraseña actual.
\end{itemize}

\subsection{Historial de Pedidos (RF-24)}
\begin{itemize}
    \item \textbf{¿Qué pasa si no tiene pedidos anteriores?} $\rightarrow$ Se muestra: <<Aún no tienes pedidos. ¡Explora nuestro catálogo!>>
\end{itemize}

\subsection{Lista de Deseos (RF-25)}
\begin{itemize}
    \item \textbf{¿Qué pasa si la lista está vacía?} $\rightarrow$ Se muestra: <<Tu lista de deseos está vacía.>>
    \item \textbf{¿Qué pasa si un producto ya no está disponible?} $\rightarrow$ Se marca como <<No disponible>> y se deshabilita agregar al carrito.
\end{itemize}

\subsection{Gestión de Productos --- Vendedor (RF-26, RF-27)}
\begin{itemize}
    \item \textbf{¿Qué pasa si no sube imágenes?} $\rightarrow$ El sistema requiere al menos una imagen para publicar.
    \item \textbf{¿Qué pasa si publica con campos vacíos?} $\rightarrow$ Se resaltan los campos obligatorios faltantes.
    \item \textbf{¿Qué pasa si da de baja un producto que está en carritos?} $\rightarrow$ Se marca como no disponible en esos carritos.
\end{itemize}

\subsection{Control de Inventario (RF-28)}
\begin{itemize}
    \item \textbf{¿Qué pasa si pone el stock en 0?} $\rightarrow$ El producto se muestra como <<Agotado>> en el catálogo.
    \item \textbf{¿Qué pasa si el stock llega a nivel bajo?} $\rightarrow$ Se muestra una alerta visual en el panel del vendedor.
\end{itemize}

\subsection{Paquetes Promocionales (RF-29)}
\begin{itemize}
    \item \textbf{¿Qué pasa si selecciona menos de 2 productos?} $\rightarrow$ El botón <<Crear paquete>> permanece deshabilitado.
    \item \textbf{¿Qué pasa si el precio del paquete es mayor que la suma individual?} $\rightarrow$ Se muestra advertencia sobre el precio.
\end{itemize}

\subsection{Pedidos Recibidos --- Vendedor (RF-30)}
\begin{itemize}
    \item \textbf{¿Qué pasa si no tiene órdenes?} $\rightarrow$ Se muestra: <<Aún no tienes pedidos. ¡Promociona tus productos!>>
    \item \textbf{¿Qué pasa si intenta regresar un pedido a un estado anterior?} $\rightarrow$ El sistema no permite retroceder el estado.
\end{itemize}

\subsection{Resumen de Ventas (RF-31)}
\begin{itemize}
    \item \textbf{¿Qué pasa si no tiene ventas?} $\rightarrow$ Se muestran los indicadores en cero y un mensaje motivacional.
\end{itemize}

\subsection{Administración --- Gestión de Usuarios (RF-32, RF-33, RF-34)}
\begin{itemize}
    \item \textbf{¿Qué pasa si intenta eliminar una cuenta?} $\rightarrow$ Se muestra un modal de confirmación: <<¿Estás seguro? Esta acción no se puede deshacer.>>
    \item \textbf{¿Qué pasa si intenta eliminarse a sí mismo?} $\rightarrow$ El sistema no permite la auto-eliminación.
\end{itemize}

\subsection{Reportes y Denuncias (RF-35)}
\begin{itemize}
    \item \textbf{¿Qué pasa si no hay reportes pendientes?} $\rightarrow$ Se muestra: <<No hay reportes pendientes por revisar.>>
    \item \textbf{¿Qué pasa si desestima un reporte?} $\rightarrow$ El reporte cambia a estado <<Resuelto>> y se archiva.
    \item \textbf{¿Qué pasa si el mismo usuario/producto recibe múltiples denuncias?} $\rightarrow$ Se agrupan y se resaltan como caso prioritario.
\end{itemize}

\newpage

% ============================================================
\section{Matriz de Navegación}
% ============================================================

La siguiente matriz indica, para cada pantalla, a qué otras pantallas puede navegar el usuario.

\subsection{Pantallas Comunes (Todos los Roles)}
\begin{longtable}{|p{5cm}|p{9cm}|}
\hline
\rowcolor{darkbg}\textcolor{white}{\textbf{Pantalla Origen}} & \textcolor{white}{\textbf{Puede ir a}} \\
\hline
Registro de Usuario & $\rightarrow$ Inicio de Sesión \\
\hline
Inicio de Sesión & $\rightarrow$ Registro de Usuario \newline $\rightarrow$ Página Principal (si Comprador) \newline $\rightarrow$ Panel del Vendedor (si Vendedor) \newline $\rightarrow$ Panel de Administración (si Admin) \\
\hline
\end{longtable}

\subsection{Pantallas del Comprador}
\begin{longtable}{|p{5cm}|p{9cm}|}
\hline
\rowcolor{darkbg}\textcolor{white}{\textbf{Pantalla Origen}} & \textcolor{white}{\textbf{Puede ir a}} \\
\hline
Página Principal / Catálogo & $\rightarrow$ Búsqueda Inteligente \newline $\rightarrow$ Detalle de Producto \newline $\rightarrow$ Carrito de Compras \newline $\rightarrow$ Lista de Deseos \newline $\rightarrow$ Mi Cuenta / Perfil \newline $\rightarrow$ Inicio de Sesión (cerrar sesión) \\
\hline
Búsqueda Inteligente & $\rightarrow$ Detalle de Producto (al seleccionar) \newline $\rightarrow$ Página Principal (resultados) \\
\hline
Detalle de Producto & $\rightarrow$ Carrito de Compras (agregar) \newline $\rightarrow$ Lista de Deseos (guardar) \newline $\rightarrow$ Página Principal (volver) \newline $\rightarrow$ Calificar Producto (sub-pantalla) \\
\hline
Carrito de Compras & $\rightarrow$ Detalle de Producto (clic en artículo) \newline $\rightarrow$ Confirmación de Pedido (checkout) \newline $\rightarrow$ Página Principal (seguir comprando) \\
\hline
Confirmación de Pedido & $\rightarrow$ Pantalla de Éxito \newline $\rightarrow$ Página Principal (seguir comprando) \newline $\rightarrow$ Historial de Pedidos \\
\hline
Mi Cuenta / Perfil & $\rightarrow$ Historial de Pedidos \newline $\rightarrow$ Lista de Deseos \newline $\rightarrow$ Página Principal \newline $\rightarrow$ Inicio de Sesión (cerrar sesión) \\
\hline
Historial de Pedidos & $\rightarrow$ Detalle de Pedido (expandir) \newline $\rightarrow$ Mi Cuenta / Perfil (volver) \\
\hline
Lista de Deseos & $\rightarrow$ Detalle de Producto \newline $\rightarrow$ Carrito de Compras (agregar) \newline $\rightarrow$ Mi Cuenta / Perfil (volver) \\
\hline
\end{longtable}

\subsection{Pantallas del Vendedor}
\begin{longtable}{|p{5cm}|p{9cm}|}
\hline
\rowcolor{darkbg}\textcolor{white}{\textbf{Pantalla Origen}} & \textcolor{white}{\textbf{Puede ir a}} \\
\hline
Panel del Vendedor (Productos) & $\rightarrow$ Formulario Nuevo Producto \newline $\rightarrow$ Editar Producto \newline $\rightarrow$ Control de Inventario \newline $\rightarrow$ Paquetes Promocionales \newline $\rightarrow$ Pedidos Recibidos \newline $\rightarrow$ Resumen de Ventas \newline $\rightarrow$ Mi Perfil \newline $\rightarrow$ Inicio de Sesión (cerrar sesión) \\
\hline
Control de Inventario & $\rightarrow$ Panel del Vendedor (volver) \\
\hline
Paquetes Promocionales & $\rightarrow$ Panel del Vendedor (volver) \\
\hline
Pedidos Recibidos & $\rightarrow$ Detalle de Orden (expandir) \newline $\rightarrow$ Resumen de Ventas \newline $\rightarrow$ Panel del Vendedor (volver) \\
\hline
Resumen de Ventas & $\rightarrow$ Pedidos Recibidos \newline $\rightarrow$ Panel del Vendedor (volver) \\
\hline
\end{longtable}

\subsection{Pantallas del Administrador}
\begin{longtable}{|p{5cm}|p{9cm}|}
\hline
\rowcolor{darkbg}\textcolor{white}{\textbf{Pantalla Origen}} & \textcolor{white}{\textbf{Puede ir a}} \\
\hline
Panel de Administración & $\rightarrow$ Gestión de Usuarios \newline $\rightarrow$ Reportes y Denuncias \newline $\rightarrow$ Mi Perfil \newline $\rightarrow$ Inicio de Sesión (cerrar sesión) \\
\hline
Gestión de Usuarios & $\rightarrow$ Detalle de Usuario (ver perfil) \newline $\rightarrow$$\rightarrow$ Modal Eliminar Cuenta \newline $\rightarrow$ Panel de Administración (volver) \\
\hline
Reportes y Denuncias & $\rightarrow$ Detalle del Reporte \newline $\rightarrow$ Gestión de Usuarios (ver reportado) \newline $\rightarrow$ Panel de Administración (volver) \\
\hline
\end{longtable}

Porcentaje de entrega 95%

\end{document}
